\documentclass[a4paper, 11pt]{article}

\usepackage[text={15cm, 23cm}]{geometry}
\usepackage[czech]{babel}
\usepackage{fontspec}
\usepackage{newcomputermodern}
\usepackage{hyperref}
\usepackage{microtype}
\usepackage{biblatex}
\usepackage{booktabs}

\newtheorem{theorem}{Věta}

\linespread{1.05}

\addbibresource{xkouma02.bib}

\begin{document}

\begin{titlepage}
    \begin{center}
        {\Huge
            \textsc{Vysoké Učení Technické v Brně}\\[0.5em]}
        {\huge
            \textsc{Fakulta Informačních Technologií}}\\
        \vspace{\stretch{0.382}}
        {\LARGE
            Teorie her -- technická zpráva\\[0.8em]
            \textbf{Induktivní důkaz Arrowova teorému s využitím SAT solveru}}\\[2em]
        \Large\today\\
        \vspace{\stretch{0.618}}
    \end{center}
    {\Large
    Bc. Ondřej Koumar\hfill\href{mailto:xkouma02@stud.fit.vutbr.cz}{xkouma02@stud.fit.vutbr.cz}}
\end{titlepage}

\section{Úvod}\label{sec:Introduction}

Cílem tohoto projektu je prakticky si vyzkoušet \emph{computer-aided proof} impossibility teorému z~teorie her, konkrétně Arrowova teorému.
Namísto klasického čistě matematického důkazu se zaměřujeme na kombinaci induktivního argumentu a~automatického ověřování bázového kroku pomocí SAT solveru.

Arrowův teorém je příkladem tvrzení, jehož důkaz vyžaduje hodně mechanické práce kvůli kombinatorické explozi problému.
Induktivní kroky lze odvodit poměrně elegantně, ale samotný bázový případ vyžaduje systematické procházení velkého prostoru možností.
Tuto mechanickou část práce zanecháme počítači.

Hlavní myšlenkou projektu je převést axiomy Arrowova teorému pro konečný počet voličů a~alternativ do formule výrokové logiky a~ověřit její splnitelnost pomocí SAT solveru.
Nesplnitelnost takové formule odpovídá neexistenci \emph{funkce společenského blahobytu} splňující všechny dané axiomy.

Nejprve stručně představíme Arrowův teorém a~použitou notaci.
Poté popíšeme induktivní krok důkazu, který umožňuje redukovat obecný případ na bázovou instanci se dvěma voliči a~třemi alternativami.
Následně se zaměříme na samotný převod tohoto bázového případu do CNF formule a~na její řešení pomocí SAT solveru.
Na závěr jsou uvedeny experimentální výsledky, včetně diskuse o~škálování a~výpočetní náročnosti celého přístupu.


\section{Formální definice Arrowova teorému}

V~této sekci zavedeme základní notaci pro voliče, alternativy a~preference.

\subsection{Základní notace}

Nechť $A$ je konečná množina \emph{alternativ} (kandidátů), kde $|A| \ge 3$.
Dále mějme množinu \emph{voličů} (agentů), kterou označíme $Q$, o~velikosti $n$ (kde $n \ge 2$).

Prostor všech možných uspořádání na množině $A$ (permutací) označíme $\mathcal{R}$.
Každý volič $i \in Q$ má svou preferenci, kterou modelujeme jako relaci na $A$.
Používáme následující značení pro vztah dvou alternativ $x, y \in A$ z~pohledu voliče $i$:
\begin{itemize}
    \item $x \succeq_i y$: Volič $i$ preferuje $x$ před $y$ nebo je mezi nimi lhostejný (slabá preference).
    \item $x \succ_i y$: Volič $i$ striktně preferuje $x$ před $y$.
\end{itemize}
My se budeme především zaměřovat na striktní preferenci.

\emph{Profil preferencí} $\rho$ je uspořádaná $n$-tice preferencí všech voličů:
\[
    \rho = (R^1, R^2, \dots, R^n) \in \mathcal{R}^n
\]
kde $R^i$ reprezentuje uspořádání $i$-tého voliče v~daném profilu.

Cílem je nalézt \emph{funkci společenského blahobytu} (Social Welfare Function) $F$.
Tato funkce agreguje preference voličů do jedné společenské preference:
\[
    F \colon \mathcal{R}^n \to \mathcal{R}
\]
Výslednou preferenci společnosti pro profil $\rho$ značíme jednoduše jako $x \succeq y$ (resp. $x \succ y$), což znamená, že společnost (podle pravidel mechanismu $F$) hodnotí $x$ alespoň tak dobře jako $y$.

\subsection{Axiomy Arrowova teorému}

Arrowův teorém zkoumá existenci funkce $F$, která by splňovala následující tři podmínky:

\begin{enumerate}
    \item \emph{Slabá Pareto optimalita}: Pokud všichni voliči striktně preferují alternativu $a$ před $b$, musí i~společnost striktně preferovat $a$ před $b$.
          Formálně:
          \[
              \forall a, b \in A \colon (\forall i \in Q \colon a \succ_i b) \implies a \succ b
          \]
    \item \emph{Nezávislost na irelevantních alternativách (IIA)}: Společenské uspořádání mezi $a$ a~$b$ závisí pouze na individuálních preferencích mezi těmito dvěma alternativami.
          Tedy, pokud voliči jednotlivě preferují $a$ nad $b$, výsledná společenská preference by s~tímto měla být konzistentní.
          Formálně, pro libovolné dva profily $\rho, \rho' \in \mathcal{R}^n$ platí:
          \begin{align*}
              \forall a, b & \in A\colon                                                                                                       \\
                           & \left[\forall i \in Q \colon (a \succeq_i b \iff a \succeq'_i b) \land (b \succeq_i a \iff b \succeq'_i a)\right] \\
                           & \implies (a \succeq b \iff a \succeq' b)
          \end{align*}
    \item \emph{Nediktátorství}: Neexistuje volič, jehož striktní preference by se vždy stala striktní preferencí společnosti, bez ohledu na názory ostatních.
          Formálně požadujeme, aby neexistoval $d \in N$ takový, že:
          \[
              \forall a, b \in A, \forall \rho \in \mathcal{R}^n \colon a \succ_d b \implies a \succ b
          \]
\end{enumerate}

\subsection{Znění teorému}

\begin{theorem}
    Pokud počet alternativ $|A| \ge 3$ a~počet voličů $|Q| \ge 2$, pak neexistuje žádná funkce společenského blahobytu $F$, která je současně tranzitivní, slabě Pareto optimální, nezávislá na irelevantních alternativách a~nediktátorská.
\end{theorem}

Jinými slovy, aby mechanismus uměl s~libovolnými preferencemi voličů (včetně remíz), byl Pareto efektivní a~nezávislý na irelevantních alternativách, musí být diktátorský.

\section{Induktivní krok důkazu}

Tradiční matematická indukce obvykle postupuje od menšího k~většímu (tj. dokážeme pro $n$ a~z~toho plyne $n+1$).
V~případě \emph{impossibility} teorémů v~teorii her je ale výhodnější postupovat opačným směrem, respektive dokázat, že pokud teorém platí pro $n$, tak po redukci platí i~pro $n-1$.
Vznikne tak implikace:
\[
    \text{platí pro } n \implies \text{platí pro } n-1
\]
a~postupnou redukcí se dostaneme až k~implikaci
\[
    \text{platí pro } n \implies \text{platí pro } k,
\]
kde $k$ je určitá konstanta.
Pro spor budeme předpokládat, že premisa platí, tím pádem musí platit i~důsledek, který bude následně vyvrácen, čímž je důkaz zkompletován.

Důkaz pak probíhá ve třech krocích:
\begin{enumerate}
    \item \emph{Redukce alternativ:} Ukážeme, že existence funkce pro $m$ alternativ implikuje existenci funkce pro $m-1$.
    \item \emph{Redukce voličů:} Ukážeme, že existence funkce pro $n$ voličů implikuje existenci funkce pro $n-1$.
    \item \emph{Bázový krok:} Pro minimální instanci ($n=2, m=3$) ověříme pomocí SAT solveru, že žádná taková funkce neexistuje.
\end{enumerate}
Pokud bázový krok selže (funkce neexistuje), pak díky implikacím z~prvních dvou kroků nutně neexistuje řešení pro žádné konečné $n$ a~$m$.

\subsection{Redukce počtu alternativ}

Předpokládejme, že existuje funkce $F_m$ pro množinu alternativ $A$ o~velikosti $m > 3$.
Naším cílem je definovat funkci $F_{m-1}$ na množině $A' = A \setminus \{x\}$, kde $x$ je libovolně zvolená alternativa, kterou z~modelu odstraníme.

Funkci $F_{m-1}$ zkonstruujeme pomocí vnoření profilů z~$A'$ do $A$.
Pro libovolný profil $\rho$ definovaný na menší množině $A'$ vytvoříme rozšířený profil $\rho^{+x}$ na množině $A$ tak, že alternativu $x$ umístíme na poslední místo v~preferencích všech voličů.
Formálně, pro každého voliče $i$ a~každou alternativu $y \in A'$ v~profilu $\rho^{+x}$ platí $y \succ_i x$.

Samotnou funkci $F_{m-1}$ pak definujeme předpisem:
\[
    F_{m-1}(\rho) = F_m(\rho^{+x})|_{A'},
\]
kde $|_{A'}$ značí zúžení preferenčního uspořádání společnosti na množinu $A'$.
Funkce $F_m$ nejprve určí společenskou preferenci na celé množině $A$, přičemž alternativa $x$ je v~individuálních preferencích zafixována jako striktně nejhorší.
Z~axiomu Pareto efektivity vyplývá, že $x$ se musí umístit na posledním místě i~ve výsledném společenském uspořádání.
Důsledkem axiomu IIA pak je, že relativní společenské pořadí alternativ uvnitř $A'$ závisí výhradně na jejich vzájemném individuálním pořadí.
Protože fixní pozice $x$ na dně žebříčků nijak nemění preference mezi ostatními prvky, dědí $F_{m-1}$ vlastnosti IIA a~nediktátorství přímo od původní funkce $F_m$.
Tímto způsobem jsme prokázali, že existence funkce splňující Arrowovy axiomy pro $m$ alternativ implikuje její existenci i~pro $m-1$ alternativ.
Tuto redukci lze aplikovat opakovaně až do dosažení bázového případu pro tři alternativy.
\subsection{Redukce počtu voličů}

Druhý indukční krok je snižování počtu voličů.
Předpokládejme existenci validní funkce $F_n$ pro množinu voličů o~velikosti $n > 2$.
Naším cílem je zkonstruovat funkci $F_{n-1}$ pro $n-1$ voličů, která si zachová vlastnosti původní funkce.

Konstrukci provedeme metodou \uv{sloučení voličů}, což budeme modelovat vynucenou shodou preferencí $n$-tého voliče s~voličem prvním.
Formálně definujeme funkci $F_{n-1}$ pro profil $\rho = (R_1, \dots, R_{n-1})$ předpisem:
\[
    F_{n-1}(R_1, \dots, R_{n-1}) = F_n(R_1, \dots, R_{n-1}, R_1)\,.
\]
Tímto jsme omezili doménu původní funkce $F_n$ na podmnožinu profilů, kde platí podmínka shody $R_1 = R_n$.
Nyní ověříme, že se vlastnosti $F_n$ přenášejí i~na novou funkci $F_{n-1}$.

\subsubsection*{Paretovská efektivita a~IIA}
Nechť $\rho = (R_1, \dots, R_{n-1})$ je libovolný profil, ve kterém všichni voliči upřednostňují $a$ před $b$, tedy určitě platí $a \succ_1 b$.
Výsledek funkce $F_{n-1}$ je naprosto shodný s~výsledkem funkce $F_n$, kde argumentem je profil $(R_1, \dots, R_{n-1}, R_1)$, tedy nutně platí $a \succ_n b$, tudíž stále platí, že $a$ je jednomyslně preferováno před $b$.
Původní funkce $F_n$ je Paretovská, proto určitě platí $a \succ b$, čímž je podmínka splněna i~pro $F_{n-1}$.

Podobně uvažujme dva profily $\rho$ a~$\sigma$, které se shodují v~preferencích na dvojici $\{a, b\}$.
Jelikož se preference $R_1$ shodují v~obou profilech, musí se shodovat i~$R_n$ v~obou profilech.
Všechny vstupy do funkce $F_n$ jsou tedy z~hlediska dvojice $\{a, b\}$ identické v~obou případech.
Díky vlastnosti IIA funkce $F_n$ je výsledná společenská preference stejná.
Tím pádem musí platit, že $F_{n-1}$ rozhoduje nezávisle na ostatních alternativách.

\subsubsection*{Nediktátorství}

Zbývá ověřit, že redukcí počtu voličů nevznikne diktátor.
Důkaz provedeme sporem.
Předpokládejme, že funkce $F_{n-1}$ má diktátora $d \in \{1,\dots,n-1\}$.
Ukážeme, že za tohoto předpokladu musí být diktátorská i funkce $F_n$.

Nechť
\[
    \rho = (R_1,\dots,R_n) \in \mathcal{R}^n
\]
je libovolný profil preferencí a~nechť volič $d$ v~tomto profilu striktně preferuje $a \succ_d b$.
Aby byl $d$ diktátorem funkce $F_n$, musí prosadit $a \succ b$ v~každém takovém profilu $\rho$.
Rozlišíme proto dva případy podle vztahu preferencí voličů $1$ a~$n$ na dvojici $\{a,b\}$, ostatní preference jsou pro oba voliče totožné.

Pokud se voliči 1 a $n$ shodují v~preferenci na dvojici $\{a, b\}$, tedy $R_1|_{\{a,b\}} = R_n|_{\{a,b\}}$, nachází se profil $\rho$ v~doméně funkce $F_{n-1} = F_n(R_1,\dots,R_{n-1},R_1)$.
Z~předpokladu, že volič $d$ je diktátorem funkce $F_{n-1}$, proto přímo plyne společenské rozhodnutí $a \succ b$ ve funkci $F_n$.

Zbývá tedy uvažovat případ, kdy se voliči 1 a~$n$ na preferenci $\{a, b\}$ neshodují, tedy $R_1|_{\{a,b\}} \ne R_n|_{\{a,b\}}$.
V~tomto případě zkonstruujeme pomocný profil
\[
    \rho' = (R_1,\dots,R_{n-1},R_n'),
\]
kde preference voliče $n$ upravíme pouze na dvojici $\{a,b\}$ tak, aby platilo $R_n'|_{\{a,b\}} = R_1|_{\{a,b\}}$ a~jeho preference vůči ostatním alternativám ponecháme beze změny.

Z~definice redukce tedy platí
\[
    F_n(\rho') = F_{n-1}(R_1,\dots,R_{n-1}),
\]
a~v~profilu $\rho'$ rozhoduje výhradně volič $d$.
Protože $d$ striktně preferuje $a \succ b$, platí $a \succ_{\rho'} b$.

Nyní porovnáme profily $\rho$ a~$\rho'$.
Tyto profily se shodují v~preferencích všech voličů kromě voliče $n$ a~navíc se shodují v~preferencích všech voličů na dvojici $\{a,b\}$ kromě právě voliče $n$.
Preference voliče $d$ na této dvojici zůstala nezměněna.
Díky axiomu IIA platí, že společenské rozhodnutí mezi $a$ a~$b$ musí být v~profilech $\rho$ a~$\rho'$ shodné.
Z~$a \succ_{\rho'} b$ tedy plyne $a \succ_{\rho} b$.

Ukázali jsme, že volič $d$ prosazuje svou striktní preferenci $a \succ b$ i~v~libovolném profilu $\rho$.
Volič $d$ je tedy diktátorem funkce $F_n$, což je spor s~předpokladem nediktátorství.

\section{Bázový krok a jeho převedení do logické formule}\label{sec:Encoding}

Nyní se zaměříme na bázový krok ($n=2, |A|=3$). Naším úkolem je převést existenci funkce $F$ na problém splnitelnosti formule výrokové logiky\footnote{Konkrétně budeme konstruovat formuli v~CNF.}.
Hledáme model, který by reprezentoval validní volební mechanismus.
Pokud SAT solver vrátí \texttt{UNSAT}, znamená to, že taková funkce neexistuje.
Transformace problému do logické formule je založena na článcích \cite{brandt2016finding} a~\cite{geist2011automated}.

\subsection{Proměnné}

Prostor všech možných vstupních profilů $\mathcal{R}^n$ je sice konečný, ale pro manuální dokazování poměrně nepraktický.
V~tomto experimentu pracujeme se striktními preferencemi.
Pro $|A|=3$ existuje $3! = 6$ možných preferencí pro jednoho voliče.
Pro $n=2$ voliče tak máme celkem $6^2 = 36$ možných vstupních profilů.

Zavedeme boolovské proměnné $w_{a, b, \rho}$, které reprezentují výstupní hodnotu funkce $F$:
\[
    w_{a, b, \rho} \equiv \text{Společnost preferuje } a \succ b \text{ v profilu } \rho.
\]
Kde $a, b \in A, a \neq b$ a~$\rho \in \mathcal{R}^n$.
Označme $p_{i, a, b, \rho}$ jako pravdivostní hodnotu výroku \uv{volič $i$ v~profilu $\rho$ striktně preferuje $a$ před $b$}.
Vstupní proměnné $p$ jsou konstanty dané konkrétním profilem $\rho$.

\subsection{Omezení}

Aby proměnné $w_{a, b, \rho}$ reprezentovaly platnou funkci společenského blahobytu, musíme vygenerovat klauzule v~CNF pro každou podmínku.

\subsubsection{Vlastnosti společenského uspořádání}
Výstupem funkce musí být lineární uspořádání. Pro každý profil $\rho \in \mathcal{R}^n$ a~všechny $a, b, c \in A$ generujeme:

\begin{enumerate}
    \item \emph{Úplnost}: Společnost musí mít názor buď $a \succ b$ nebo $b \succ a$.
          \[
              w_{a, b, \rho} \lor w_{b, a, \rho}
          \]
    \item \emph{Asymetrie}: Společnost nemůže preferovat $a \succ b$ a~současně $b \succ a$.
          \[
              \neg w_{a, b, \rho} \lor \neg w_{b, a, \rho}
          \]
    \item \emph{Tranzitivita}: Pokud $a$ je preferováno před $b$ a~$b$ preferováno před $v$, pak musí být $a$ preferováno před $c$.
          \[
              (\neg w_{a, b, \rho} \lor \neg w_{b, c, \rho} \lor w_{a, c, \rho})
          \]
          Tato klauzule odpovídá implikaci $(w_{a, b, \rho} \land w_{b, c, \rho}) \implies w_{a, c, \rho}$.
\end{enumerate}

\subsubsection{Slabá Pareto optimalita}
Pokud v~profilu $\rho$ všichni voliči upřednostňují $a$ před $b$, musí tak činit i~společnost.
Pokud platí $\forall i \colon p_{i, a, b, \rho}$ (což zjistíme už při generování), přidáme klauzuli:
\[
    w_{a, b, \rho}\,.
\]

\subsubsection{Nezávislost na irelevantních alternativách}
Toto je výpočetně nejnáročnější podmínka, protože počet klauzulí roste kvadraticky s~počtem profilů.
Procházíme všechny dvojice profilů $\rho, \rho' \in \mathcal{R}^n$ a~všechny dvojice alternativ $a, b \in A$.
Pokud se preference všech voličů vůči dvojici $\{a, b\}$ v~profilech $\rho$ a~$\rho'$ shodují, musí se shodovat i~výstupní hodnota $w$.
Definujme si pomocný predikát $\text{Match}(\rho, \rho', a, b)$.
\[
    \text{Match}(\rho, \rho', a, b) \equiv \bigwedge_{i \in Q} (p_{i, a, b, \rho} \leftrightarrow p_{i, a, b, \rho'})
\]
Pokud je $\text{Match}$ pravdivý, přidáme na výstup:
\[
    (w_{a, b, \rho} \lor \neg w_{a, b, \rho'}) \land (\neg w_{a, b, \rho} \lor w_{a, b, \rho'}),
\]
což odpovídá ekvivalenci
\[
    w_{a, b, \rho} \iff w_{a, b, \rho'}\,.
\]

\subsubsection{Nediktátorství}
Chceme dokázat, že neexistuje diktátor. Pro každého voliče $i$ tedy musí existovat alespoň jeden případ, profil $\rho$ a~dvojice $a, b$, kdy se společenská preference liší od jeho preference.
To zapíšeme jako jednu velkou disjunkci pro každého voliče $i$:
\[
    \bigvee_{\rho \in \mathcal{R}^n, a, b \in A} (p_{i, a, b, \rho} \land \neg w_{a, b, \rho})\,.
\]

\section{Experimenty}\label{sec:Experiments}


Pro ověření bázového kroku induktivního důkazu byl implementován skript v~jazyce Python, jehož úkolem je automaticky generovat odpovídající instanci CNF formule.
Ta je následně řešena pomocí SAT solveru \textsc{Minisat}.
Skript pro zadaný počet voličů a~alternativ nejprve vygeneruje všechny možné profily striktních preferencí.
Na základě těchto profilů jsou poté postupně konstruovány klauzule odpovídající vlastnostem relace slabé preference a~jednotlivým axiomům Arrowova teorému.

Hlavním cílem experimentální části je praktická verifikace Arrowova teorému pro bázový případ $N=2$, $M=3$.
Nicméně skript byl implementován tak, aby se jednotlivé axiomy mohly odebírat, čímž vzniklo zajímavé okno pro zkoumání velikosti a~jiných vlastností výsledné formule, k~čemuž se stručně vyjádříme na konci této sekce.

Pro každou instanci danou počtem voličů $N$ a~alternativ $M$ jsou generovány axiomy Pareto optimality, nezávislosti na irelevantních alternativách a~nediktátorství.
Solver následně rozhodne, zda existuje model reprezentující funkci společenského blahobytu splňující danou množinu axiomů.

Experimenty byly prováděny ve dvou základních režimech:
\begin{enumerate}
    \item \emph{Plná formulace teorému}, ve které jsou zahrnuty všechny axiomy Arrowova teorému.
    \item \emph{Relaxovaná formulace}, ve které je vždy vynechán právě jeden axiom a je testováno, zda se formule stane splnitelnou.
\end{enumerate}

Základní experiment odpovídá instanci s~$N=2$ voliči a~$M=3$ alternativami.
Tento případ, spolu s~induktivním krokem, postačuje k~důkazu Arrowova teorému v~obecnosti (\uv{up to bugs in the solver} -- a~v~našem případě i~v~generátoru formule -- nicméně pro všechny testované instance byl výsledek očekávaný).
Výsledky pro větší hodnoty $N$ a~$M$ jsou proto chápány především jako ilustrativní a~slouží k~posouzení výpočetní náročnosti celé metody.

\subsection{Výsledky experimentů}

Výsledky experimentů pro bázový případ $N=2$, $M=3$ jsou shrnuty v~tabulce~\ref{tab:basic-results}.
\begin{table}[h]
    \centering
    \begin{tabular}{lrrrrl}
        \toprule
        Chybějící axiom & Profily & Klauzule & Generace [s] & Řešení [s] & Výsledek \\
        \midrule
        žádný           & 36      & 2432     & 0.0015       & 0.0001     & UNSAT    \\
        Pareto          & 36      & 2378     & 0.0013       & 0.0000     & SAT      \\
        IIA             & 36      & 704      & 0.0006       & 0.0000     & SAT      \\
        Nedikt.         & 36      & 2430     & 0.0012       & 0.0001     & SAT      \\
        \bottomrule
    \end{tabular}
    \caption{Výsledky pro $N=2$, $M=3$. }
    \label{tab:basic-results}
\end{table}

V~případě plné formulace teorému solver vrací \texttt{UNSAT}, což znamená, že neexistuje žádná funkce společenského blahobytu splňující všechny axiomy současně.
Tento výsledek odpovídá tvrzení Arrowova teorému.

Po odebrání jednotlivých axiomů se formule stává splnitelnou.
V~případě vynechání podmínky nediktátorství nalezený model odpovídá diktátorskému mechanismu, kdy společenská preference kopíruje preference jednoho z~voličů.
To potvrzuje, že diktatura je jediným mechanismem splňujícím zbývající axiomy.

Při odebrání podmínky nezávislosti na irelevantních alternativách by solver teoreticky měl nalézt model odpovídající systému \emph{Borda count} -- nezkoumáno podrobně.
Dále při odstranění této podmínky dochází k~výraznému snížení počtu klauzulí, což naznačuje, že právě IIA představuje výpočetně nejnáročnější část.


\subsection{Škálování a výpočetní náročnost}

Ačkoliv je pro samotný důkaz Arrowova teorému postačující bázový případ $N=2$, $M=3$, implementovaný skript umožňuje generovat a~řešit i~větší instance.
Vybrané výsledky pro různé kombinace počtu voličů a~alternativ jsou uvedeny v~tabulce~\ref{tab:scaling-results}.

\begin{table}[h]
    \centering
    \begin{tabular}{rrrrrrl}
        \toprule
        $N$ & $M$ & Profily & Klauzule     & Generace [s] & Řešení [s] & Výsledek \\
        \midrule
        3   & 3   & 216     & 37\,749      & 0.017        & 0.0004     & UNSAT    \\
        4   & 3   & 1296    & 645\,898     & 0.276        & 0.0075     & UNSAT    \\
        5   & 3   & 7776    & 11\,432\,183 & 4.720        & 0.125      & UNSAT    \\
        2   & 4   & 576     & 1\,024\,706  & 0.431        & 0.0058     & UNSAT    \\
        \bottomrule
    \end{tabular}
    \caption{Vybrané experimenty pro větší instance.}
    \label{tab:scaling-results}
\end{table}

Z~výsledků je patrné, že počet klauzulí roste velmi rychle, zejména s~rostoucím počtem alternativ.
Navzdory tomu zůstává doba samotného řešení SAT solverem relativně nízká.
Dominantní část výpočetního času připadá na konstrukci CNF formule, nikoliv na její řešení.
Nutno dodat, že skript byl během experimentování několikrát optimalizován a~jeho finální podoba je zhruba $80\times$ rychlejší než původní verze.

Experimentální výsledky tak potvrzují nejen platnost Arrowova teorému v~bázovém případě, ale také demonstrují praktickou použitelnost SAT solverů pro verifikaci podobných nemožnostních tvrzení v~teorii společenské volby.


\section{Závěr}\label{sec:Conclusion}

V~rámci projektu jsme ukázali, že Arrowův teorém lze prakticky ověřit kombinací induktivního argumentu a~automatického řešení bázového kroku pomocí SAT solveru.
Induktivní redukce umožňuje zmenšit problém na instanci, kterou lze mechanicky převést do CNF formule a~tu následně ověřit.

Experimenty potvrdily, že pro bázový případ se dvěma voliči a~třemi alternativami je výsledná formule nesplnitelná, což odpovídá tvrzení Arrowova teorému.
Zároveň se ukázalo, že ačkoliv velikost formule roste velmi rychle s~počtem voličů a~alternativ, samotné řešení SAT solverem zůstává poměrně efektivní.
Hlavním úzkým místem je generování formule, nikoliv její řešení.

Použitý přístup představuje elegantní způsob, jak induktivně dokazovat teorémy, jejichž bázové kroky jsou příliš rozsáhlé pro manuální práci.
Zároveň však vyvstává otázka, do jaké míry lze takové důkazy považovat za korektní.
V~praxi se tento typ důkazu v~odborné literatuře běžně používá, zejména u~tvrzení jako Arrowův teorém.

Striktně vzato však nelze s~absolutní jistotou tvrdit, že celý řetězec nástrojů -- od generátoru formule přes SAT solver až po samotné běhové prostředí -- je zcela bezchybný.
Pokud tuto předpokládanou korektnost přijmeme, lze metodu považovat za důkaz.
V~opačném případě se musíme spokojit s~tvrzením, že výsledek je \uv{dost pravděpodobně správný}.

\pagebreak

\nocite{*}
\printbibliography

\end{document}